\documentclass[11pt]{article}
\usepackage[margin=1in]{geometry}
\usepackage{amsmath,amssymb,amsthm,booktabs,hyperref}
\newtheorem{theorem}{Theorem}
\newtheorem{lemma}{Lemma}
\newtheorem{proposition}{Proposition}
\theoremstyle{remark}
\newtheorem{remark}{Remark}
\title{A smaller generic grid in the unit-distance argument}
\author{AI-generated research draft prepared with Codex and Claude}
\date{October 10, 2026}
\begin{document}
\maketitle
\paragraph{Scope and attribution.}
This is an independent follow-up to
\href{https://github.com/openai/math/tree/fd4aeeb2ee4fc729c18d98444fed42fd0529eeeb/preprints/A-power-saving-for-planar-unit-distances-September-23-2026}{OpenAI's original manuscript},
pinned at commit \texttt{fd4aeeb2ee4fc729c18d98444fed42fd0529eeeb}.
It is not an official OpenAI publication. Mathematical priority and
independent formal verification are not claimed. Exact checks supplement
the proofs below and do not verify the imported source lemmas.

\paragraph{Summary of the bounds.}
The source selects a $20$-by-$5$ grid, or $100$ edges. The first
version of this note reduced the witness to $9$-by-$4$ in general and
$6$-by-$3$ in the common-line branch. The present version gives
\[
 \begin{array}{lcc}
 \toprule
 \text{branch} & \text{grid} & \text{edges}\\
 \midrule
 \text{general (rank two)} & 8\times4 & 32\\
 \text{one locus is a line (rank two)} & 6\times4 & 24\\
 \text{both loci the same line} & 4\times3 & 12\\
 \bottomrule
 \end{array}
\]
All three are conditional on the source's generic-grid and
private-valuation lemmas. A single $8$-by-$4$ generic grid contains the
smaller sub-grids, so one selection covers every case.

\section{The symmetry group of the differential system}
\label{sec:symmetries}
Retain the original differential system, with $J(x,y)=(-y,x)$:
\begin{equation}\label{eq:system}
 2(p_i-r_\ell)\cdot(H_i-U_\ell)
   =R_{i\ell}(A_i-A'_\ell),\qquad A_i-A'_\ell\ne0 .
\end{equation}
Here $H_i=\partial p_i$, $U_\ell=\partial r_\ell$,
$A_i=\partial g_i/g_i$ and $A'_\ell=\partial h_\ell/h_\ell$ for the
source's $k_0$-derivation $\partial$ and its normalized ratios
$g_i,h_\ell$. Since $h_1=\lambda_{11}/\lambda_{11}=1$, one has
$A'_1=0$ automatically.

For $T\in\mathcal U^2$ and $\omega,c\in\mathcal U$ with $c\ne0$, the
replacement
\begin{equation}\label{eq:gauge}
 H_i\mapsto c\bigl(H_i-T-\omega Jp_i\bigr),\quad
 U_\ell\mapsto c\bigl(U_\ell-T-\omega Jr_\ell\bigr),\quad
 A_i\mapsto cA_i,\quad A'_\ell\mapsto cA'_\ell
\end{equation}
preserves \eqref{eq:system}, including the nonvanishing of
$A_i-A'_\ell$. The isometric part preserves each left side because
$x\cdot Jx=0$, and the common factor $c$ multiplies both sides. The
scaled data are the values of the derivation $c\,\partial$; the shifted
data need not be values of a new derivation, and none is needed: the
later linear algebra uses only the preserved equations.

\paragraph{Normalization.}
Let $v=r_2-r_1$. The anchors are distinct real points at finite indices,
so $v\cdot v\ne0$ in the ultraproduct field. Apply \eqref{eq:gauge} with
$T=U_1$, $c=1$ and
\[
 \omega=\frac{(U_2-U_1)\cdot Jv}{v\cdot v}.
\]
The new data satisfy $\widetilde U_1=0$ and
$\widetilde U_2\cdot Jv=0$, hence $\widetilde U_2=\tau v$ with
\[
 \tau=\frac{(U_2-U_1)\cdot v}{v\cdot v}\in F .
\]
If $\tau\ne0$, apply \eqref{eq:gauge} again with $T=0$, $\omega=0$,
$c=1/\tau$. In both cases the final normalized data satisfy
\begin{equation}\label{eq:normalized}
 \widetilde U_1=0,\qquad
 \widetilde U_2\in\{0,v\},\qquad
 \widetilde A'_1=0 .
\end{equation}
Since $v\in B_0(r_1,r_2)^2$, the normalized right-hand data require,
after all anchor positions have been adjoined, at most
\begin{equation}\label{eq:cost}
 2(j-2)+(j-1)=3j-5
\end{equation}
additional field generators: two for each of
$\widetilde U_3,\ldots,\widetilde U_j$ and one for each of
$\widetilde A'_2,\ldots,\widetilde A'_j$.

\begin{proposition}[Reduction of the final algebraic witness]
\label{prop:general}
Conditional on the source's generic-grid and private-valuation lemmas,
the final algebraic contradiction can be run with an $8$-by-$4$ grid
instead of a $20$-by-$5$ grid.
\end{proposition}
\begin{proof}
Use the source's generic-grid lemma with $N=8$, $j=4$. Let $K$ be the
algebraic closure, inside the source's algebraically closed ambient
field, of the field generated by $B_0$, the four anchor positions, and
the normalized data \eqref{eq:normalized}. This does not adjoin the
negligible-height field. By \eqref{eq:cost} the data cost at most seven
generators. If the left locus has dimension $d_1\in\{1,2\}$, the left
points jointly retain transcendence degree at least $8d_1-7$ over $K$.
If none were individually generic, subadditivity would give at most
$8(d_1-1)=8d_1-8$, a contradiction. A generic left point $p$ therefore
survives over $K$; write $M=K(p)$.

The source constructs private odd valuations directly after enlargement
of the coefficient field, using only the original joint genericity of
the positions over $B_0$. That proof applies over this $K$. Hence the
four radicands have independent square classes in $M$, and the Kummer
statement of the source holds in $M(R_1,\ldots,R_4)$.

If the rows $p-r_1,p-r_2,p-r_3$ have rank two, the first three
equations of \eqref{eq:system}, with the normalized data on the right,
have coefficients and right sides in $M(R_1,R_2,R_3)$ and a determinant
that is a nonzero $M$-linear combination of $R_1,R_2,R_3$; Cramer's rule
places $\widetilde H_x,\widetilde H_y,\widetilde A$ in that field. The
automorphism changing only the sign of $R_4$ fixes these unknowns.
Applied to the fourth equation and subtracted, it gives
$2R_4(\widetilde A-\widetilde A'_4)=0$, so
$\widetilde A=\widetilde A'_4$, contradicting \eqref{eq:system}.
The rank condition fails only when both loci are the same line, which
is handled with fewer points in Section~\ref{sec:common-line}; the
$8$-by-$4$ grid contains the $4$-by-$3$ sub-grid used there.
\end{proof}

Eight is the least value this count allows: with $N=7$ the two sides of
$7d_1-7>7(d_1-1)$ are equal.

\section{Collinear anchors need six left points}
\label{sec:collinear}
The proposition's hypotheses are symmetric under exchanging the two
vertex sets. Suppose that one of the two loci is a line, and name the
sides so that the anchor locus $V_2$ is that line. Write
\[
 r_\ell=r_1+\sigma_\ell v,\qquad
 \sigma_\ell=\frac{(r_\ell-r_1)\cdot v}{v\cdot v}\in F,\qquad
 \sigma_1=0,\ \sigma_2=1 .
\]
The derivation satisfies the Leibniz rule, so
\[
 U_\ell-U_1=\partial(\sigma_\ell v)
   =\sigma_\ell\,\partial v+(\partial\sigma_\ell)\,v .
\]
With $\omega$ as in Section~\ref{sec:symmetries}, the vector
$\partial v-\omega Jv$ is orthogonal to $Jv$, hence equals
$\kappa v$ with $\kappa=(\partial v\cdot v)/(v\cdot v)$, because
$v\cdot Jv=0$ and $v\cdot v\ne0$. Therefore every normalized anchor
velocity is parallel to the line:
\begin{equation}\label{eq:parallel}
 \widetilde U_\ell=U_\ell-U_1-\omega J(r_\ell-r_1)
   =\mu_\ell v,\qquad
 \mu_\ell=\sigma_\ell\kappa+\partial\sigma_\ell,\qquad
 \mu_1=0,\ \mu_2=\kappa=\tau .
\end{equation}
After the scaling step, $\mu_2\in\{0,1\}$. The normalized data are then
the scalars $\mu_3,\ldots,\mu_j$ and
$\widetilde A'_2,\ldots,\widetilde A'_j$, so they cost at most
\begin{equation}\label{eq:collinear-cost}
 (j-2)+(j-1)=2j-3
\end{equation}
generators over $B_0(r_1,\ldots,r_j)$.

\begin{proposition}[One locus is a line]
\label{prop:collinear}
Conditional on the same source lemmas, if either locus is a line the
final contradiction can be run with a $6$-by-$4$ grid whose four
anchors lie on that line.
\end{proposition}
\begin{proof}
Take $N=6$, $j=4$ with the anchors on the line. By
\eqref{eq:collinear-cost} the data cost at most five generators, and
$6d_1-5>6(d_1-1)$ leaves a generic left point $p$ over $K$. If
$V_1=V_2$, Section~\ref{sec:common-line} applies to a sub-grid.
Otherwise $p$ is generic on the irreducible locus $V_1\ne V_2$, so it
does not lie on the anchor line, and the rows $p-r_1,p-r_2,p-r_3$ have
rank two. The right sides of \eqref{eq:system} are
$2\mu_\ell(p-r_\ell)\cdot v-R_\ell\widetilde A'_\ell\in M(R_\ell)$, and
the Cramer and sign-change steps of
Proposition~\ref{prop:general} apply verbatim.
\end{proof}

\section{The common-line branch needs only twelve selected edges}
\label{sec:common-line}
Suppose both loci are the same line. Then $p-r_\ell=s_\ell v$ with
$s_\ell\in M^\times$, only $h=v\cdot\widetilde H$ and $\widetilde A$
enter \eqref{eq:system}, and the right sides are
$2s_\ell\mu_\ell(v\cdot v)-R_\ell\widetilde A'_\ell$. Three anchors
suffice: the first two equations determine $h$ and $\widetilde A$, and
the third radical supplies the sign change. For $j=3$ the collinear
cost \eqref{eq:collinear-cost} is three generators,
$\mu_3,\widetilde A'_2,\widetilde A'_3$. Select $N=4$ left points. In
this case $d_1=1$, and
\[
 4\cdot1-3>4\cdot0 ,
\]
so a left point survives as generic over $K$. The first two equations
have coefficient matrix
\[
 \begin{pmatrix}2s_1&-R_1\\2s_2&-R_2\end{pmatrix}
\]
with determinant $2(s_2R_1-s_1R_2)$, nonzero by the independent
single-radical basis and $s_1,s_2\ne0$. Cramer's rule places $h$ and
$\widetilde A$ in the field of the first two radicals. Flipping only
$R_3$ in the third equation and subtracting forces
$\widetilde A=\widetilde A'_3$, contrary to the preserved nonzero
scalar difference. Thus a $4$-by-$3$ witness, or twelve edges, suffices
for this branch. Three left points would give $3\cdot1-3=0$ and prove
nothing.

\section{What the method cannot do}
\label{sec:limits}
The reductions above are exhausted at these numbers for the following
reasons.

\begin{remark}[The symmetry group is five-dimensional]
The transformations \eqref{eq:gauge} form a group of dimension five:
two translations, one rotation, the scalar shift
$A_i\mapsto A_i+\beta$, $A'_\ell\mapsto A'_\ell+\beta$, and the
scaling. The shift is already spent on $A'_1=0$. Dilations
$H_i\mapsto H_i+\mu p_i$, $U_\ell\mapsto U_\ell+\mu r_\ell$ change the
left side of \eqref{eq:system} by $2\mu L_{i\ell}$, which is not of
the form $R_{i\ell}(\beta_i-\beta'_\ell)$, so no further continuous
symmetry is available. Hence $3j-5$ is the least generator count
obtainable by normalizing the anchor data alone.
\end{remark}

\begin{remark}[Four anchors are needed in the rank-two case]
Three equations in the three unknowns $H_x,H_y,A$ are solved by
Cramer's rule and leave no equation for the sign change. The count
$N>\text{cost}$ is also sharp for the subadditivity argument.
\end{remark}

\begin{remark}[Choosing the derivation]
The source's derivation may be chosen to vanish on a rational subspace
of the anchor differentials that avoids all $d\lambda_{i\ell}$, which
removes anchor data in favorable cases. This gives no unconditional
saving: all normalized data may be algebraic over $k_0(t)$ for a
single $t$, in which case the relations have coefficients outside the
coefficient field. We therefore do not use it.
\end{remark}

None of this provides an explicit power saving $\delta$ for the
unit-distance bound: the generic-grid extraction and negligible-height
passage remain asymptotic, and no quantitative replacement is
established here.

\paragraph{Reproducibility.}
The accompanying README documents \texttt{python -B run\_checks.py}.
The scripts, source manifest, and test expectations are in this repository.
\end{document}
